\documentclass[11pt]{article}
\usepackage[margin=1in]{geometry}
\usepackage{amsmath,amssymb,amsthm}
\usepackage{enumitem}

\newtheorem{theorem}{Theorem}
\newtheorem{proposition}[theorem]{Proposition}
\newtheorem{corollary}[theorem]{Corollary}
\newtheorem{remark}[theorem]{Remark}

\title{d5\_240: Reverse-search reduction for the remaining graph-side target}
\author{}
\date{2026-03-22}

\begin{document}
\maketitle

\section{Purpose}
After the exhaustion of the surfaced mode-switch family, the honest graph-side
route is to search from the smallest reduced target signature rather than from a
wider local selector grammar.
This note isolates that target at the reduced permutation level.
The outcome is sharper than a heuristic: inside the exact diagonal/anti-diagonal
one-row adjacent-transposition family, the \emph{only} base mechanism that
produces a single orbit is the omitted-row defect, and once that base is fixed,
a one-point cocycle defect already yields a single full grouped orbit for every
odd modulus.

\section{The reduced family}
Fix an odd modulus $m>2$.
For parameters
\[
 a\in\{+1,-1\},\qquad b,d\in\mathbb Z_m,
\]
write
\[
 g(s)=as+b.
\]
Define a map
\[
 T_{a,b,d,h}\colon \mathbb Z_m^2\to\mathbb Z_m^2,
 \qquad (s,u)\mapsto (s+1,\tau_s(u)),
\]
where for $s\neq d$ the row permutation $\tau_s$ is the adjacent transposition
swapping $g(s)$ and $g(s)+1$, and on the defect row $s=d$ one uses either
\begin{enumerate}[label=\textup{(\roman*)},leftmargin=2.5em]
\item the identity map (the \emph{omitted-row defect}), written $h=\infty$, or
\item the adjacent transposition swapping $h$ and $h+1$ for some $h\in\mathbb Z_m$.
\end{enumerate}
So this is exactly the moving adjacent-transposition family with one defect row.

\section{Affine conjugacy to a canonical model}
\begin{proposition}[canonical reduction]
\label{prop:canonical-reduction}
For every $a,b,d$ the affine change of variables
\[
 \Phi_{a,b,d}(s,u)=\bigl(a(s-d),\ u-(ad+b)\bigr)
\]
conjugates $T_{a,b,d,h}$ to a canonical map $T_{+,0,0,h'}$ with defect row $0$,
where $h'=h-(ad+b)$ if $h\in\mathbb Z_m$, and $h'=\infty$ in the omitted-row
case.
\end{proposition}

\begin{proof}
The first coordinate sends the defect row $d$ to $0$.
If $u=g(s)=as+b$, then under $\Phi_{a,b,d}$ one gets
\[
 u' = as+b-(ad+b)=a(s-d)=s'.
\]
So the active adjacent pair $\{g(s),g(s)+1\}$ is sent to $\{s',s'+1\}$.
The same computation identifies the defect-row transposition with the canonical
pair $\{h',h'+1\}$.
Thus the whole family is conjugate to the canonical defect-row-$0$ model.
\end{proof}

Hence it suffices to study the canonical family on $\mathbb Z_m^2$ with defect
row $0$ and baseline row transpositions
\[
 \tau_s=(s\ \ s+1)\qquad (s=1,\dots,m-1),
\]
where addition in the symbol list is modulo $m$.

\section{The omitted-row defect is the unique single-orbit base mechanism}
\begin{proposition}[omitted-row first return]
\label{prop:omit-first-return}
In the canonical omitted-row case $T_{+,0,0,\infty}$, the first-return map to
row $0$ is the translation
\[
 R(u)=u-1\qquad (u\in\mathbb Z_m).
\]
Therefore $T_{+,0,0,\infty}$ has a single orbit of size $m^2$.
\end{proposition}

\begin{proof}
One full return from row $0$ applies, in order, the adjacent transpositions
\[
 (1\ \ 2),\ (2\ \ 3),\ \dots,\ (m-2\ \ m-1),\ (m-1\ \ 0).
\]
So the first-return permutation is
\[
 R=(m-1\ \ 0)\cdots(2\ \ 3)(1\ \ 2).
\]
A direct evaluation gives $R(u)=u-1$ for every $u\in\mathbb Z_m$.
Since $R$ is an $m$-cycle on row $0$ and the first coordinate advances by one at
each step, the whole map on $\mathbb Z_m^2$ is a single orbit of size $m^2$.
\end{proof}

\begin{proposition}[shifted defect first return]
\label{prop:shift-first-return}
In the canonical shifted-defect case $T_{+,0,0,h}$ with $h\in\mathbb Z_m$, the
first-return map to row $0$ is
\[
 R_h = R\circ(h\ \ h+1),
\]
where $R(u)=u-1$ is the omitted-row return map.
It has cycle structure
\[
 [1,m-1].
\]
Hence $T_{+,0,0,h}$ has orbit decomposition
\[
 [m,\ m(m-1)]
\]
and in particular never has a single orbit.
\end{proposition}

\begin{proof}
The defect-row transposition $(h\ \ h+1)$ is applied first, and then the same
baseline return $R$ as in Proposition~\ref{prop:omit-first-return}.
So the return map is exactly $R_h=R\circ(h\ \ h+1)$.
By translation in $u$ it suffices to treat $h=0$.
Then
\[
 R_0 = R\circ(0\ \ 1),
\]
so $0$ is fixed, while
\[
 1\mapsto m-1\mapsto m-2\mapsto\cdots\mapsto 1
\]
forms one cycle of length $m-1$.
The full map therefore has one orbit above the fixed return point and one orbit
above the $(m-1)$-cycle, giving orbit sizes $m$ and $m(m-1)$.
\end{proof}

\begin{theorem}[exact reduced candidate reduction]
\label{thm:reduced-candidate-reduction}
Inside the exact diagonal/anti-diagonal one-row adjacent-transposition family
$T_{a,b,d,h}$, the following are equivalent:
\begin{enumerate}[label=\textup{(\arabic*)},leftmargin=2.8em]
\item $T_{a,b,d,h}$ has a single orbit on $\mathbb Z_m^2$;
\item the defect row is omitted, i.e. $h=\infty$.
\end{enumerate}
So the omitted-row diagonal/anti-diagonal family is the unique single-orbit
base target in this exact reduced class.
\end{theorem}

\begin{proof}
By Proposition~\ref{prop:canonical-reduction}, everything reduces to the
canonical cases $T_{+,0,0,\infty}$ and $T_{+,0,0,h}$.
The omitted-row case is single-orbit by
Proposition~\ref{prop:omit-first-return}.
Every shifted-defect case splits into two orbits by
Proposition~\ref{prop:shift-first-return}.
\end{proof}

\section{The smallest full grouped target}
The reduced full grouped model is the skew product on
$\mathbb Z_m^2\times\mathbb Z_m$:
\[
 \widehat T(x,v)=\bigl(T(x),\ v+\psi(x)\bigr).
\]
If $T$ has one base orbit $\mathcal O$ of length $m^2$, then the orbit structure
of $\widehat T$ is controlled by the total cocycle sum on one base cycle.

\begin{proposition}[single-cycle cocycle criterion]
\label{prop:cocycle-criterion}
Let $T$ be a permutation of a finite set $X$ with one orbit $\mathcal O$ of
length $N$.
Let $\psi\colon X\to\mathbb Z_m$, and define
\[
 \widehat T(x,v)=\bigl(T(x),\ v+\psi(x)\bigr).
\]
If the cycle sum
\[
 c=\sum_{x\in\mathcal O}\psi(x)
\]
is invertible modulo $m$, then $\widehat T$ has one orbit of size $Nm$.
\end{proposition}

\begin{proof}
After $N$ steps the base coordinate returns and the fiber coordinate advances by
$c$.
If $c$ is invertible modulo $m$, repeated returns generate the whole fiber
$\mathbb Z_m$.
So every point on the unique base orbit lifts to one orbit of size $Nm$.
\end{proof}

\begin{corollary}[point and two-point defects]
\label{cor:point-and-two-point}
For the omitted-row base target $T_{a,b,d,\infty}$ on odd $m>2$:
\begin{enumerate}[label=\textup{(\arabic*)},leftmargin=2.8em]
\item a one-point cocycle defect on one omitted-edge endpoint gives a single
      full orbit of size $m^3$;
\item a two-point cocycle defect on both omitted-edge endpoints also gives a
      single full orbit of size $m^3$.
\end{enumerate}
\end{corollary}

\begin{proof}
By affine conjugacy, reduce to the canonical omitted-row base.
Then the omitted-edge endpoints are $(0,0)$ and $(0,1)$.
A one-point defect has cycle sum $1$, and a two-point defect has cycle sum $2$.
For odd $m$, both are invertible in $\mathbb Z_m$.
Apply Proposition~\ref{prop:cocycle-criterion} together with
Theorem~\ref{thm:reduced-candidate-reduction}.
\end{proof}

\section{What this does and does not close}
This note still does \emph{not} realize the remaining graph-side theorem $G1$
locally.
What it does close is the reduced reverse-search question:
\begin{quote}
Inside the smallest exact diagonal/anti-diagonal one-row defect family, the
only honest base target is the omitted-row adjacent-transposition model, and the
smallest full grouped target is that base together with a one-point omitted-edge
cocycle defect.
\end{quote}
So any remaining local search should be judged against that exact signature,
not against a wider generic selector grammar.

\section{Exact search sanity check}
The companion script
\texttt{d5\_239\_G1\_reduced\_omit\_target\_search.py}
performs the exact finite search over all parameters
\[
 a\in\{\pm1\},\quad b\in\mathbb Z_m,\quad d\in\mathbb Z_m,
\quad h\in\{\infty\}\cup\mathbb Z_m
\]
for checked odd moduli
\[
 m=3,5,7,9,11,13,15,17,19.
\]
It confirms:
\begin{enumerate}[label=\textup{(\arabic*)},leftmargin=2.8em]
\item omission count $=2m^2$ exactly;
\item shifted-defect single-orbit count $=0$ exactly;
\item canonical omitted-row return cycle size $=[m]$;
\item canonical shifted-defect return cycle size $=[1,m-1]$;
\item canonical one-point and two-point cocycle defects both give one full
      orbit of size $m^3$.
\end{enumerate}
This is only a sanity check; the main statements above are symbolic.

\end{document}
